\subsection{PASSAGE FROM LIE ALGEBRA MORPHISMS TO LIE GROUP MORPHISMS}

Lemma 1. Let \(\mathbf{G}\) be a Lie group germ and \(\mathfrak{h}\) a Lie subalgebra of \(\mathbf{L}(\mathbf{G})\) admitting a topological supplement. The union of the \(\mathfrak{gh}\) (resp. \(\mathfrak{hg}\) ) for \(g\in\mathbf{G}\) is an integrable vector subbundle of \(\mathbf{T}(\mathbf{G})\) .

By considering the left trivialization of \(\mathrm{T}(\mathrm{G})\) (§ 2, no. 3), it is seen immediately that the \(g\mathfrak{h}\) , for \(g\in \mathrm{G}\) , are the fibres of a vector subbundle \(\mathrm{E}\) of \(\mathrm{T}(\mathrm{G})\) . Let \(g\in \mathrm{G}\) . The set of \((\mathrm{L}_a)_g\) , where \(a\in \mathfrak{h}\) , is equal to \(g\mathfrak{h}\) . Now, if \(a\) and \(b\) belong to \(\mathfrak{h}\) , then \([\mathrm{L}_a,\mathrm{L}_b] = \mathrm{L}_{[a,b)}\) and \([a,b]\in \mathfrak{h}\) . Hence \(\mathrm{E}\) is integrable (Differentiable and Analytic Manifolds, R, 9.3.3 (iv)). The argument is similar for the \(\mathfrak{h}g\) .

The integral foliation (Differentiable and Analytic Manifolds, R, 9.3.2) of the union of the \(\mathfrak{gh}\) (resp. \(\mathfrak{hg}\) ) is called the left (resp. right) foliation of G associated with \(\mathfrak{h}\) .

THEOREM 1. Let \(\mathbf{G}\) and \(\mathbf{H}\) be Lie group germs and \(f\) a continuous morphism of \(\mathbf{L}(\mathbf{G})\) into \(\mathbf{L}(\mathbf{H})\) .

\begin{enumerate}
\def\labelenumi{(\roman{enumi})}
\item
  There exist an open Lie subgroup germ \(\mathbf{G}'\) of \(\mathbf{G}\) and a morphism \(\phi\) of \(\mathbf{G}'\) into \(\mathbf{H}\) such that \(f = \mathbf{L}(\phi)\) .
\item
  Let \(\mathrm{G}_1, \mathrm{G}_2\) be open Lie subgroup germs of G and \(\phi_i\) a morphism of \(\mathrm{G}_i\) into H such that \(f = \mathrm{L}(\phi_i)\) for \(i = 1, 2\) . Then \(\phi_1\) and \(\phi_2\) coincide on a neighbourhood of \(e\) . Let \(p_1: \mathrm{G} \times \mathrm{H} \to \mathrm{G}\) , \(p_2: \mathrm{G} \times \mathrm{H} \to \mathrm{H}\) be the canonical projections. For all \((g,h) \in \mathrm{G} \times \mathrm{H}\) , let \(f_{g,h}\) be the mapping \(ga \mapsto hf(a)\) of \(\mathrm{T_g(G)} = g\mathrm{L(G)}\) into \(\mathrm{T_h(H)} = h\mathrm{L(H)}\) . By considering the left trivializations of \(\mathrm{T(G)}\) and \(\mathrm{T(H)}\) , it is seen immediately that the \(f_{g,h}\) define a morphism of \(p_1^*\mathrm{T(G)}\) into \(p_{2}^{*}\mathbf{T}(\mathrm{H})\) . Let \(a\) be the graph of \(f\) ; it is a closed Lie subalgebra of \(\mathbf{L}(\mathrm{G}) \times \mathbf{L}(\mathrm{H})\) which admits \(\{0\} \times \mathbf{L}(\mathrm{H})\) as topological supplement. For all \((g,h) \in \mathrm{G} \times \mathrm{H}\) , the graph of \(f_{g,h}\) is \((g,h)\) . \(a\) . The union of these graphs is an integrable vector subbundle of \(\mathrm{T}(\mathrm{G} \times \mathrm{H})\) (Lemma 1). Then there exist (Differentiable and Analytic Manifolds, R, 9.3.7) an open neighbourhood \(\mathrm{U}\) of \(e_{\mathrm{G}}\) in \(\mathrm{G}\) and an analytic mapping \(\phi\) of \(\mathrm{U}\) into \(\mathrm{H}\) such that \(\phi(e_{\mathrm{G}}) = e_{\mathrm{H}}\) and \(\mathrm{T}_{g}(\phi) = f_{g,\phi(g)}\) for all \(g \in \mathrm{U}\) . In particular, \(\mathrm{T}_{e_{\mathrm{G}}}(\phi) = f\) .
\end{enumerate}

Let \(\mathbf{V}\) be an open neighbourhood of \(e_{\mathrm{G}}\) in \(\mathbf{G}\) such that, for \((s,t)\in \mathbf{V}\times \mathbf{V}\) , the products \(st\) and \(\phi (s)\phi (t)\) are defined and \(st\in \mathbf{U}\) . Consider the mappings \(\alpha_{1},\alpha_{2}\) of \(\mathbf{V}\times \mathbf{V}\) into \(\mathbf{H}\) defined by

\[
\alpha_ {1} (s, t) = \phi (t s), \quad \alpha_ {2} (s, t) = \phi (t) \phi (s).
\]

Then \(\alpha_{1}(t,e) = \phi (t) = \alpha_{2}(t,e)\) . On the other hand, let \(t\) be fixed in V and \(\beta_{i}\) be the mapping \(s\mapsto \alpha_i(s,t)\) of V into H. Then, for all \(s\in \mathrm{V}\) and all \(a\in \mathrm{L}(\mathrm{G})\) ,

\[
\begin{array}{r l} \mathrm{T} _ {s} (\beta_ {1}) (s a) & = \mathrm{T} _ {t s} (\phi) (t s a) = f _ {t s, \phi (t s)} (t s a) \\ & = \phi (t s) f (a) = f _ {s, \beta_ {1} (s)} (s a) \\ \mathrm{T} _ {s} (\beta_ {2}) (s a) & = \phi (t) \mathrm{T} _ {s} (\phi) (s a) = \phi (t) f _ {s, \phi (s)} (s a) \\ & = \phi (t) \phi (s) f (a) = f _ {s, \beta_ {2} (s)} (s a). \end{array}
\]

Hence (Differentiable and Analytic Manifolds, R, 9.3.7) \(\alpha_{1}\) and \(\alpha_{2}\) coincide on a neighbourhood of \((e_{\mathrm{G}}, e_{\mathrm{G}})\) . The restriction of \(\phi\) to a sufficiently small symmetric open neighbourhood of \(e_{\mathrm{G}}\) is therefore a morphism of Lie group germs, whence (i).

Let \(\mathrm{G}_1, \mathrm{G}_2, \phi_1, \phi_2\) be as in (ii) and let us prove that \(\phi_1, \phi_2\) coincide on a neighbourhood of \(e_{\mathrm{G}}\) . There exists an open neighbourhood W of \(e_{\mathrm{G}}\) such that \(\phi_1(t s) = \phi_1(t)\phi_1(s)\) , \(\phi_2(t s) = \phi_2(t)\phi_2(s)\) for all \(s\) , \(t\) in W. Then, if \(s \in W\) and \(a \in L(G)\) ,

\[
\mathrm{T} _ {s} \left(\phi_ {i}\right) (s a) = \phi_ {i} (s) \mathrm{T} _ {e} \left(\phi_ {i}\right) (a) = \phi_ {i} (s) f (a) = f _ {s, \phi _ {i} (s)} (s a)
\]

for \(i = 1,2\) . As \(\phi_1(e_{\mathrm{G}}) = e_{\mathrm{H}} = \phi_2(e_{\mathrm{G}})\) , it follows from Differentiable and Analytic Manifolds, R, 9.3.7 that \(\phi_1\) and \(\phi_2\) coincide on a neighbourhood of \(e_{\mathrm{G}}\) .

COROLLARY 1. Let G and H be two Lie group germs. If L(G) and L(H) are isomorphic, G and H are locally isomorphic.

This follows from Theorem 1 and § 1, no. 10, Proposition 21.

COROLLARY 2. Let \(\mathbf{G}\) be a Lie group germ. If \(\mathrm{L}(\mathbf{G})\) is commutative, \(\mathbf{G}\) is locally isomorphic to the additive Lie group \(\mathrm{L}(\mathbf{G})\) .

The Lie algebra of the additive group \(\mathbf{L}(\mathbf{G})\) is isomorphic to \(\mathbf{L}(\mathbf{G})\) . Hence it suffices to apply Corollary 1.

COROLLARY 3. Let G be a Lie group. If L(G) is commutative, G contains a commutative open subgroup.

There exists an open Lie subgroup germ U of G which is commutative (Corollary 2). Let V be a neighbourhood of e such that \(V^{2} \subset U\) . Then xy = yx for all x, y in V. Hence the subgroup of G generated by V is commutative; it is obviously open.

